\documentclass[lualatex,ja=standard]{bxjsarticle}
\usepackage{amsmath,amssymb}
\usepackage{enumitem}
\setlist[enumerate]{itemsep=0.6em, topsep=0.5em}
\pagestyle{empty}

\begin{document}

\noindent\textbf{数IA 思考演習（標準） 解答}
\hrule
\vspace{0.8em}

\begin{enumerate}
\item
\[
-7\le 2x-5\le 7
\Rightarrow -2\le 2x\le 12
\Rightarrow -1\le x\le 6
\]
$x$ は整数より
\[
x=-1,0,1,2,3,4,5,6
\]

\item 異なる2実数解の条件は判別式 $D>0$。
\[
D=(a+1)^2-4(a-2)=a^2-2a+9=(a-1)^2+8
\]
これは常に正なので、すべての実数 $a$ で条件を満たす。

\item
\begin{enumerate}[label=(\arabic*)]
\item
\[
y=-(x^2-4x)+1=-(x-2)^2+5
\]
より頂点は $(2,5)$。
\item
\[
-x^2+4x+1=0
\Leftrightarrow x^2-4x-1=0
\]
\[
x=\frac{4\pm\sqrt{16+4}}{2}=2\pm\sqrt5
\]
よって共有点は
\[
(2-\sqrt5,0),\ (2+\sqrt5,0)
\]
\end{enumerate}

\item
\[
(x+y)^2=x^2+2xy+y^2
\]
より
\[
25=13+2xy \Rightarrow xy=6
\]
したがって $x,y$ は
\[
t^2-5t+6=0
\]
の2解なので
\[
t=2,3
\]
ゆえに
\[
(x,y)=(2,3),(3,2)
\]

\item
\begin{enumerate}[label=(\arabic*)]
\item 余弦定理より
\[
12^2=10^2+10^2-2\cdot10\cdot10\cos A
\]
\[
144=200-200\cos A
\Rightarrow \cos A=\frac{7}{25}
\]
よって
\[
\angle A=\cos^{-1}\!\left(\frac{7}{25}\right)
\]
\item
\[
S=\frac12\cdot10\cdot10\sin A=50\sin A
\]
\[
\sin A=\sqrt{1-\cos^2A}
=\sqrt{1-\left(\frac{7}{25}\right)^2}
=\frac{24}{25}
\]
より
\[
S=50\cdot\frac{24}{25}=48
\]
\end{enumerate}

\item $0^\circ<\theta<180^\circ$ で $\tan\theta<0$ より第2象限。
\[
\tan\theta=\frac{\sin\theta}{\cos\theta}=-\frac43
\]
比より $\sin:\cos=4:-3$、また
\[
\sin^2\theta+\cos^2\theta=1
\]
から
\[
\sin\theta=\frac45,\quad \cos\theta=-\frac35
\]

\item 全事象は
\[
\binom92=36
\]
和が偶数になるのは「奇数2枚」または「偶数2枚」。
奇数は5枚、偶数は4枚なので
\[
\binom52+\binom42=10+6=16
\]
よって確率は
\[
\frac{16}{36}=\frac49
\]

\item 全体は
\[
\binom93=84
\]
3色すべて異なるには赤1・白1・黒1の選び方なので
\[
\binom41\binom31\binom21=4\cdot3\cdot2=24
\]
よって
\[
24\ \text{通り}
\]

\item
\begin{itemize}
\item[(ア)] 中央値70点より、20番目と21番目の値の平均が70。よって少なくとも20人が70点以上とは限らない（例：20番目69点、21番目71点）。
\item[(イ)] 同様に必ずしも成り立たない。
\item[(ウ)] 平均68点なので、全員が68点未満だと平均も68点未満になる。したがって最高点は68点以上。
\end{itemize}
よって必ず正しいのは
\[
\text{(ウ)}
\]

\item 定義域より
\[
x+5\ge0,\ x-1\ge0 \Rightarrow x\ge1
\]
両辺を2乗して
\[
x+5=(x-1)^2=x^2-2x+1
\]
\[
x^2-3x-4=0
\Rightarrow (x-4)(x+1)=0
\]
\[
x=4,-1
\]
定義域 $x\ge1$ より
\[
x=4
\]
\end{enumerate}

\end{document}
